% BHU PYQ -- Professional Exam Template
\documentclass[11pt,a4paper]{article}

\usepackage[a4paper,top=2.5cm,bottom=2.5cm,left=2.8cm,right=2.8cm,headheight=14pt]{geometry}
\usepackage{amsmath,amssymb}
\usepackage[T1]{fontenc}
\usepackage[utf8]{inputenc}
\usepackage{lmodern}
\usepackage{microtype}
\usepackage{fancyhdr}
\usepackage{enumitem}
\usepackage{listings}
\usepackage{hyperref}
\usepackage{lastpage}
\usepackage{parskip}

\hypersetup{colorlinks=true,urlcolor=black,linkcolor=black,
  pdftitle={BVCA-101: Introduction To Computers -- BHU PYQ}}

\pagestyle{fancy}
\fancyhf{}
\renewcommand{\headrulewidth}{0.4pt}
\renewcommand{\footrulewidth}{0.4pt}
\fancyhead[L]{\small\textit{Banaras Hindu University}}
\fancyhead[R]{\small\textit{BVCA-101: Introduction To Computers}}
\fancyfoot[C]{\small Page \thepage\ of \pageref{LastPage}}

\lstset{
  basicstyle=\small\ttfamily,
  frame=single,
  breaklines=true,
  columns=flexible,
  keepspaces=true
}

\newlist{parts}{enumerate}{2}
\setlist[parts,1]{label=(\alph*),leftmargin=*,itemsep=4pt,topsep=3pt}
\setlist[parts,2]{label=(\roman*),leftmargin=*,itemsep=2pt,topsep=2pt}
\newcommand{\pts}[1]{\hfill\textit{[#1]}}

\begin{document}

\begin{center}
  {\large\bfseries BANARAS HINDU UNIVERSITY}\\[2pt]
  {\normalsize B.Voc. Semester I Examination, 2022-23}\\[6pt]
  \rule{0.8\linewidth}{0.5pt}\\[6pt]
  {\large\bfseries Paper: BVCA-101 --- Introduction To Computers}\\[4pt]
  {\normalsize Time: Three Hours \hfill Maximum Marks: 70}
\end{center}

\rule{\linewidth}{0.4pt}

\smallskip
\noindent\textit{Instructions: Attempt all questions. Questions of Section A, B and C carry 15, 5 and 1 marks each respectively.}

\bigskip


\medskip\noindent\textbf{Question 1.} Discuss the generations of computers based on their
characteristics and processing speed.
\centerline{\textbf{OR}}
Based on block diagram of computer discuss its
different components.

\medskip\noindent\textbf{Question 2.} What is Compiler ? Describe how compiler differs
from interpreter ?
\centerline{\textbf{OR}}
Define operating system and give a detailed account
on different types of operating system.

\bigskip\noindent{\large\bfseries SECTION --- B | )}
\rule{\linewidth}{0.4pt}\smallskip


\medskip\noindent\textbf{Question 3.} What is DOS ? Describe the internal and external
commands of DOS.

\medskip\noindent\textbf{Question 4.} Provide the differences between the procedural and
object oriented programming. |

\medskip\noindent\textbf{Question 5.} Write short notes of the following :
\begin{parts}
  \item Overloading
  \begin{parts}
    \item Polymorphism
  \end{parts}
\end{parts}

\medskip\noindent\textbf{Question 6.} Define mail merge in MS-Office.

\medskip\noindent\textbf{Question 7.} Whatare the advantages and disadvantages of WAN ?

\medskip\noindent\textbf{Question 8.} What are the differences between primary | and
secondary memory ?

\medskip\noindent\textbf{Question 9.} Explain the different types of 4 common network
topologies.

\medskip\noindent\textbf{Question 10.} What is formatting in Excel ? Explain the formatting
features available in Excel for the numeric data.

\bigskip\noindent{\large\bfseries SECTION --- C}
\rule{\linewidth}{0.4pt}\smallskip

Note: Answer in briefest possible way. |

\medskip\noindent\textbf{Question 11.} Who is the father of Computers ?
\begin{parts}
  \item James Gosling
  \item Charles Babbage
  \item Dennis Ritchie
  \item Bjarne Stroustrup
\end{parts}

\medskip\noindent\textbf{Question 12.} Which of the following is the correct abbreviation of COMPUTER ?
\begin{parts}
  \item Commonly Occupied Machines Used in
  Technical and Educational Research
  \item Commonly Operated Machines Used in
  Technical and Environmental Research
  \item Commonly Oriented Machines Used in
  Technical and Educational Research
  \item Commonly Operated Machines Used in
  | Technical and Educational Research
\end{parts}

\medskip\noindent\textbf{Question 13.} What is the full form of CPU ?
\begin{parts}
  \item Computer Processing Unit
  \item Computer Principle Unit
  \item Central Processing Unit
  \item Control Processing Unit
\end{parts}

\medskip\noindent\textbf{Question 14.} Which of the following computer language is written in binary codes only ?
\begin{parts}
  \item Pascal
  \item machine language
  \item C
  \item C\#
\end{parts}

\medskip\noindent\textbf{Question 15.} Which of the following is the smallest unit of data in a computer ?
\begin{parts}
  \item Bit
  \item KB
  \item Nibble
  \item Byte
\end{parts}

\medskip\noindent\textbf{Question 16.} Which of the following service allows a user to log in to another computer somewhere on the Internet ?
\begin{parts}
  \item e-mail
  \item UseNet
  \item Telnet
  \item FIP
\end{parts}

\medskip\noindent\textbf{Question 17.} The software designed to perform a specific task:
\begin{parts}
  \item Synchronous Software
  \item Package Software
  \item Application Software
  \item System Software
\end{parts}

\medskip\noindent\textbf{Question 18.} What do you call a specific instruction designed to do a task ?
\begin{parts}
  \item Command
  \item Process
  \item Task
  \item Instruction
\end{parts}

\medskip\noindent\textbf{Question 19.} Public domain software is usually :
\begin{parts}
  \item System supported
  \item Source supported
  \item Community supported
  \item Programmer supported
\end{parts}

\medskip\noindent\textbf{Question 20.} OSS stands for :
\begin{parts}
  \item Open System Service
  \item Open Source Software
  \item Open System Software
  \item Open Synchronized Software
\end{parts}

\vfill
\rule{\linewidth}{0.4pt}
\begin{center}\small
\href{https://bhu-exam.vercel.app/}{Click here to visit our website}\\[4pt]\href{https://me-aryan.vercel.app/}{Click here to visit our contributor}\\[4pt]\href{https://quashan-me.vercel.app/}{Click here to visit our contributor}
\end{center}

\end{document}
